% sudodoc-toc.tex — Elegant TOC and sectioning system
% Consistent with sudodoc-titlepage design language
% Compile: lualatex sudodoc-toc.tex (run TWICE for TOC + overlays)
\documentclass[12pt]{article}

% ─── Packages ───
\usepackage{fontspec}
\usepackage{geometry}
\usepackage{xcolor}
\usepackage{tikz}
\usepackage{graphicx}
\usepackage{titlesec}
\usepackage{titletoc}
\usepackage{parskip}
\usepackage{calc}
\usepackage{fancyhdr}
\usepackage{lastpage}

% ─── Fonts ───
\setsansfont{Latin Modern Sans}
\setmainfont{Latin Modern Roman}
\setmonofont{Latin Modern Mono}

% ─── Colors (matching title page) ───
\definecolor{titleDark}{HTML}{1A1A2E}
\definecolor{titleAccent}{HTML}{16213E}
\definecolor{titleRule}{HTML}{0F3460}
\definecolor{titleLight}{HTML}{E8E8E8}
\definecolor{titleMid}{HTML}{888888}

% ─── Geometry ───
\geometry{
  paper=a4paper,
  top=30mm,
  bottom=30mm,
  left=30mm,
  right=25mm,
  headheight=15pt,
}

% ══════════════════════════════════════════════════════════
% SECTIONING: Chapter, Section, Subsection styling
% ══════════════════════════════════════════════════════════

% ── Chapter-style headings (using \section as "chapters") ──
\titleformat{\section}
  {\normalfont\Large\bfseries\color{titleDark}}
  {}{0pt}
  {\vspace{6pt}%
   \begin{tikzpicture}[baseline=(title.base)]
     \node[anchor=west, inner sep=0pt] (title) {\insertsection};
     \draw[titleRule, line width=1.2pt]
       ([yshift=-3pt]title.south west) -- ([yshift=-3pt]title.south east);
   \end{tikzpicture}%
  }
  [\vspace{2pt}]

\titleformat{\subsection}
  {\normalfont\large\bfseries\color{titleRule}}
  {}{0pt}
  {\hspace{4pt}%
   \tikz[baseline=(num.base)]{
     \node[fill=titleRule, text=white, rounded corners=2pt,
           inner xsep=6pt, inner ysep=2pt,
           font=\footnotesize\bfseries] (num)
       {\thesubsection};
   }%
   \hspace{6pt}\insertsubsection%
  }
  [\vspace{2pt}]

\titleformat{\subsubsection}
  {\normalfont\normalsize\bfseries\color{titleMid}}
  {}{0pt}
  {\hspace{14pt}{\color{titleRule}\textbullet}\hspace{8pt}\insertsubsubsection}
  [\vspace{1pt}]

% ── Section spacing ──
\titlespacing{\section}{0pt}{24pt plus 4pt minus 2pt}{12pt plus 2pt minus 1pt}
\titlespacing{\subsection}{0pt}{18pt plus 3pt minus 1pt}{8pt plus 1pt minus 1pt}
\titlespacing{\subsubsection}{0pt}{12pt plus 2pt minus 1pt}{6pt}

% ══════════════════════════════════════════════════════════
% TOC: Table of Contents styling
% ══════════════════════════════════════════════════════════

% ── TOC Title ──
\renewcommand{\contentsname}{\vspace*{-10pt}%
  {\fontsize{24}{28}\selectfont\bfseries\color{titleDark}Contents}}

% ── Section entries in TOC ──
\titlecontents{section}[0pt]
  {\addvspace{14pt}\fontsize{12}{14}\selectfont\bfseries\color{titleDark}}
  {\thecontentslabel\hspace{10pt}}
  {}
  {\hspace{1em}%
   \tikz[baseline=-0.5ex]{
     \draw[titleRule!40, line width=0.5pt] (0,0) -- (3cm,0);
     \node[anchor=west, font=\fontsize{10}{12}\selectfont\color{titleMid}]
       at (3.2cm,0) {\thecontentspage};
   }%
  }
  \addvspace{4pt}

% ── Subsection entries in TOC ──
\titlecontents{subsection}[16pt]
  {\addvspace{6pt}\fontsize{11}{13}\selectfont\color{titleAccent}}
  {\thecontentslabel\hspace{8pt}}
  {}
  {\hspace{1em}%
   \tikz[baseline=-0.5ex]{
     \draw[titleMid!40, line width=0.4pt] (0,0) -- (2.5cm,0);
     \node[anchor=west, font=\fontsize{10}{12}\selectfont\color{titleMid}]
       at (2.7cm,0) {\thecontentspage};
   }%
  }
  \addvspace{2pt}

% ── Subsubsection entries in TOC ──
\titlecontents{subsubsection}[32pt]
  {\addvspace{3pt}\fontsize{10}{12}\selectfont\color{titleMid}}
  {}
  {}
  {\hspace{1em}%
   \tikz[baseline=-0.5ex]{
     \draw[titleMid!25, line width=0.3pt] (0,0) -- (2cm,0);
     \node[anchor=west, font=\fontsize{9}{11}\selectfont\color{titleMid}]
       at (2.2cm,0) {\thecontentspage};
   }%
  }

% ══════════════════════════════════════════════════════════
% CHAPTER PAGE: Decorative full-page opener
% ══════════════════════════════════════════════════════════
\newcommand{\chapterpage}[2]{%
  % #1 = Chapter number
  % #2 = Chapter title
  \clearpage
  \thispagestyle{empty}
  \begin{tikzpicture}[remember picture, overlay]
    % ── Full page dark background ──
    \fill[titleDark]
      (current page.north west) rectangle (current page.south east);

    % ── Large chapter number watermark ──
    \node[anchor=south east,
          font=\fontsize{160}{160}\selectfont\bfseries,
          text=white, opacity=0.06]
      at ([xshift=-20pt, yshift=20pt]current page.south east) {#1};

    % ── Accent rule at ~45% ──
    \draw[titleRule, line width=0.4pt]
      ([xshift=40pt, yshift=-0.42\paperheight]current page.north west)
      -- ([xshift=-40pt, yshift=-0.42\paperheight]current page.north east);

    % ── "Chapter N" label ──
    \node[anchor=south west, text width=0.7\paperwidth]
      at ([xshift=0.14\paperwidth, yshift=0.06\paperheight]current page.center)
      {{\fontsize{11}{13}\selectfont\color{titleMid}%
        CHAPTER\hspace{8pt}\fontsize{13}{15}\selectfont\color{titleLight}#1}};

    % ── Chapter title ──
    \node[anchor=north west, text width=0.7\paperwidth, align=flush left]
      at ([xshift=0.14\paperwidth, yshift=0.08\paperheight]current page.center)
      {{\fontsize{28}{34}\selectfont\bfseries\color{white}#2}};

    % ── Thin decorative line under title ──
    \draw[titleRule, line width=1.2pt]
      ([xshift=0.14\paperwidth, yshift=0.005\paperheight]current page.center)
      -- ([xshift=0.55\paperwidth, yshift=0.005\paperheight]current page.center);

  \end{tikzpicture}
  \clearpage
}

% ══════════════════════════════════════════════════════════
% HEADERS / FOOTERS
% ══════════════════════════════════════════════════════════
\pagestyle{fancy}
\fancyhf{}
\fancyhead[L]{%
  {\fontsize{8}{10}\selectfont\color{titleMid}\leftmark}}
\fancyhead[R]{%
  {\fontsize{8}{10}\selectfont\color{titleMid}%
    A Comprehensive Study of Distributed Systems}}
\fancyfoot[C]{%
  {\fontsize{8}{10}\selectfont\color{titleMid}%
    \rule{2cm}{0.3pt}\hspace{6pt}\thepage\hspace{6pt}\rule{2cm}{0.3pt}}}
\renewcommand{\headrulewidth}{0pt}
\renewcommand{\footrulewidth}{0pt}

% ─── Begin document ───
\begin{document}

% ══════════════════════════════════════════════════════════
% TOC PAGE
% ══════════════════════════════════════════════════════════
\thispagestyle{empty}
\begin{tikzpicture}[remember picture, overlay]
  % ── Top accent bar (matching light title page) ──
  \fill[titleDark]
    (current page.north west)
    rectangle ([yshift=-6pt]current page.north east);
  % ── Left accent stripe ──
  \fill[titleRule]
    ([xshift=50pt]current page.north west)
    rectangle ([xshift=54pt]current page.south west);
\end{tikzpicture}

\vspace{18pt}
\hspace{22pt}\tableofcontents

% ══════════════════════════════════════════════════════════
% CHAPTER 1
% ══════════════════════════════════════════════════════════
\chapterpage{1}{Foundations of Distributed Systems}

\section{Introduction}
Distributed systems form the backbone of modern computing infrastructure.
From cloud services to social networks, the ability to coordinate multiple
autonomous machines is fundamental to scalability and reliability.

\subsection{Motivation and Goals}
The primary motivation for studying distributed systems lies in their
ubiquity. Every time you use a web application, stream a video, or check
your email, you are interacting with a distributed system.

\subsubsection{Why Distributed?}
Centralized systems, while simpler to reason about, face inherent
limitations in scale, fault tolerance, and geographic distribution.
A single point of failure can bring down an entire service.

\subsubsection{Key Challenges}
Consistency, availability, and partition tolerance form the well-known
CAP theorem triangle, forcing system designers to make deliberate
trade-offs.

\subsection{Historical Context}
The field has evolved from early remote procedure calls in the 1980s
through the rise of the internet, to modern microservices architectures
and serverless computing paradigms.

\section{System Models}
Formal models provide the theoretical foundation for reasoning about
distributed systems. These models abstract away implementation details
to focus on essential properties.

\subsection{Synchronous vs.\ Asynchronous}
In synchronous models, message delivery is bounded by a known delay.
Asynchronous models make no such guarantees, making them more realistic
but harder to reason about.

\subsection{Failure Models}
Understanding different types of failures---crash, omission, Byzantine---is
essential for designing fault-tolerant systems that can withstand
adversarial conditions.

\subsubsection{Crash Failures}
A crash failure occurs when a process halts prematurely and never
recovers. This is the simplest failure model and forms the basis for
most practical fault-tolerance mechanisms.

\subsubsection{Byzantine Failures}
Byzantine failures are the most general and most difficult to handle.
A Byzantine process may send contradictory information to different
parts of the system, making consensus especially challenging.

% ══════════════════════════════════════════════════════════
% CHAPTER 2
% ══════════════════════════════════════════════════════════
\chapterpage{2}{Consensus Algorithms}

\section{The Consensus Problem}
At the heart of every distributed system lies the challenge of getting
independent nodes to agree on a single value or course of action, despite
communication delays and potential failures.

\subsection{FLP Impossibility}
The seminal result by Fischer, Lynch, and Paterson (1985) proved that
deterministic consensus is impossible in purely asynchronous systems
with even a single failure. This impossibility result has shaped decades
of research.

\subsection{Paxos}
Leslie Lamport's Paxos protocol is one of the most influential consensus
algorithms. Despite its reputation for being difficult to understand,
its core ideas are surprisingly elegant.

\subsubsection{Basic Paxos}
Basic Paxos guarantees consensus among a set of processes, provided a
majority are responsive. It operates in two phases: prepare and accept,
using proposal numbers to ensure safety.

\subsubsection{Multi-Paxos}
Multi-Paxos optimizes the basic protocol for sequences of decisions by
electing a stable leader, eliminating the prepare phase for subsequent
values and significantly improving throughput.

\section{Raft}
Raft was designed as an understandable alternative to Paxos, decomposing
the consensus problem into independent subproblems: leader election,
log replication, and safety.

\subsection{Leader Election}
Raft uses randomized timeouts to trigger leader elections. When a follower
does not receive a heartbeat within its election timeout, it becomes a
candidate and requests votes from its peers.

\subsection{Log Replication}
Once a leader is established, clients send commands to the leader, which
appends them to its log and replicates them across followers using
AppendEntries RPCs.

% ══════════════════════════════════════════════════════════
% CHAPTER 3
% ══════════════════════════════════════════════════════════
\chapterpage{3}{Fault Tolerance in Practice}

\section{Replication Strategies}
Replication is the primary mechanism for achieving fault tolerance in
distributed systems. Different strategies offer varying trade-offs
between consistency, performance, and complexity.

\subsection{Leader-Based Replication}
In leader-based replication, all writes go through a single leader node
that propagates changes to followers. This simplifies consistency but
introduces a single point of coordination.

\subsection{Leaderless Replication}
Leaderless (or quorum-based) replication allows any node to accept writes,
using quorums to ensure consistency. Dynamo-style systems use this
approach for high availability.

\section{Failure Detection}
Detecting failures accurately and quickly is crucial for maintaining
system health. However, distinguishing between a failed process and a
temporarily slow network is notoriously difficult.

\subsection{Heartbeat Mechanisms}
Most practical systems use heartbeat-based failure detection, where
nodes periodically exchange messages. Absence of a heartbeat within
a timeout triggers a failure suspicion.

\subsection{Phi Accrual Failure Detector}
The phi accrual failure detector improves upon simple timeouts by
maintaining a probabilistic model of network delays, producing a
continuous suspicion level rather than a binary alive/dead verdict.

\end{document}
